\documentclass[11pt]{article}
\usepackage{amsmath}
\usepackage{geometry}
\geometry{margin=1in}

\title{Circuit Simulation and Breadboard Lab}
\author{}
\date{}

\begin{document}

\maketitle

\section*{Part 1: Simulated Circuits}

\subsection*{1. Resistors in Series}

Two $100\,\Omega$ resistors were connected in series with a $10\,V$ source.

Total resistance:
\[
R_{total} = 100 + 100 = 200\,\Omega
\]

Current:
\[
I = \frac{V}{R} = \frac{10}{200} = 0.05\,A
\]

Voltage across each resistor:
\[
V = IR = (0.05)(100) = 5\,V
\]

Node voltages:
\begin{itemize}
\item Source node: $10\,V$
\item Middle node: $5\,V$
\item Ground node: $0\,V$
\end{itemize}

Adding a third $100\,\Omega$ resistor:

\[
R_{total} = 300\,\Omega
\]

\[
I = \frac{10}{300} = 0.0333\,A
\]

Voltage across each resistor:
\[
V = (0.0333)(100) = 3.33\,V
\]

Node voltages:
\begin{itemize}
\item $10\,V$
\item $6.67\,V$
\item $3.33\,V$
\item $0\,V$
\end{itemize}

Adding resistance in series increases total resistance and reduces current while dividing the voltage across more nodes.

\subsection*{2. Resistors in Parallel}

Two $100\,\Omega$ resistors were connected in parallel across $10\,V$.

Voltage across each resistor:
\[
V = 10\,V
\]

Current through each resistor:
\[
I = \frac{10}{100} = 0.1\,A
\]

Total current:
\[
I_{total} = 0.1 + 0.1 = 0.2\,A
\]

Adding a third $100\,\Omega$ resistor:

\[
I_3 = 0.1\,A
\]

\[
I_{total} = 0.3\,A
\]

Voltage remains constant across branches while current increases when resistors are added in parallel.

\subsection*{3. Series and Parallel Circuit}

Two $100\,\Omega$ resistors were placed in parallel with a third $100\,\Omega$ resistor in series.

Parallel equivalent:
\[
R_p = \frac{100 \times 100}{100 + 100} = 50\,\Omega
\]

Total resistance:
\[
R_{total} = 100 + 50 = 150\,\Omega
\]

Total current:
\[
I = \frac{10}{150} = 0.0667\,A
\]

Voltage across series resistor:
\[
V = (0.0667)(100) = 6.67\,V
\]

Voltage across parallel branch:
\[
10 - 6.67 = 3.33\,V
\]

Current in each parallel resistor:
\[
I = \frac{3.33}{100} = 0.0333\,A
\]

This circuit behaves as a combination of voltage division (series) and current division (parallel).

\subsection*{4. RC Circuit}

Given:
\[
R = 100\,\Omega
\]
\[
C = 10\,\mu F
\]

Time constant:

\[
\tau = RC
\]

\[
\tau = (100)(10\times10^{-6})
\]

\[
\tau = 0.001\,s = 1\,ms
\]

Charging behavior:

\[
V_c = V(1-e^{-t/\tau})
\]

Discharging behavior:

\[
V_c = V_0 e^{-t/\tau}
\]

After one time constant:

\[
V = 0.63V_{final} \quad \text{(charging)}
\]

\[
V = 0.37V_{initial} \quad \text{(discharging)}
\]

\subsection*{5. RC Low-Pass Filter}

Given:

\[
R = 1\,k\Omega
\]

\[
C = 1\,\mu F
\]

Time constant:

\[
\tau = RC = (1000)(1\times10^{-6}) = 0.001\,s
\]

Cutoff frequency:

\[
f_c = \frac{1}{2\pi RC}
\]

\[
f_c \approx 159\,Hz
\]

For input frequencies near or above $159\,Hz$, the output waveform becomes smoother and lower in amplitude because the capacitor prevents rapid voltage changes.

\section*{Part 2: Real Circuits}

\subsection*{6. Breadboard Setup}

Current through a $1\,\Omega$ resistor with $5\,V$ supply:

\[
I = \frac{5}{1} = 5\,A
\]

The Raspberry Pi power supply can provide approximately $2\,A$, so this current would exceed the supply limit and could cause overheating or shutdown.

Using a $10\,k\Omega$ resistor:

\[
I = \frac{5}{10000} = 0.0005\,A
\]

\[
I = 10\,mA
\]

Higher resistance limits current and prevents excessive heat.

Measured voltage across resistor:

\[
V_{measured} = 5.222\,V
\]

Expected voltage:

\[
\approx 5\,V
\]

\subsection*{7. Real RC Circuit}

Given:

\[
R = 10\,k\Omega
\]

\[
C = 100\,\mu F
\]

Time constant:

\[
\tau = RC
\]

\[
\tau = (10000)(100\times10^{-6})
\]

\[
\tau = 1\,s
\]

After one time constant during discharge:

\[
V = 0.33(5.2) = 1.716\,V
\]

Measured time to reach (faulty equipment) $\approx1.716\,V$:

\[
t = 1\,s
\]

Differences between theoretical and measured values may arise from component tolerances, measurement uncertainty, and supply voltage variations.

\end{document}